\clearpage
\properlane{abundance-amid-need}{The feast sustains those who still hunger}
The Lord prepares a feast for all peoples; the Shepherd's guest lacks nothing; seekers of the Lord are promised good. Paul, however, still knows hunger. The apparent contradiction is inside the appointments themselves. It cannot be resolved by pretending that a hungry Christian has simply failed to believe, or by declaring bodily need beneath spiritual attention. God's present provision includes strength in adversity and the assistance of other people, while the feast's final abundance surpasses everything that can be consumed now.

\subsection{The counterexample of the wealthy wrongdoer}
Augustine makes the difficulty concrete in his exposition of Psalm 33. A believer sees a rich wrongdoer grow old in comfort and receive an expensive funeral. Meanwhile the believer seeks God and remains poor. If the psalm promised an observable redistribution of wealth in every lifetime, the scene would appear to refute it. The temptation then is practical as well as intellectual: imitation of the dishonest seems to offer the good that piety has failed to obtain.\footnote{Augustine, \work{Enarrationes} 33, second exposition, 13.}

Augustine's answer changes the measure of wealth. Justice, truth, charity, faith and endurance fill the heart with goods that cannot be purchased in the market. Christ himself is the living bread. A person can possess an ivory bed and lack this bread; another can lack money yet be rich in the life God gives.\footnote{Augustine, \work{Enarrationes} 33, second exposition, 14.} The distinction is not between something real called money and an unreal consolation called religion. It is between finite possessions and goods that determine whether a human life is rightly ordered and finally fulfilled.

That answer does not require the sufferer to report that hunger feels pleasant. It requires that the pain of lack not decide the worth of fidelity. Augustine knows the inducement to fraud: someone says that business cannot be conducted honestly, or that fear of God will leave nothing to eat. The psalm's promise challenges precisely that bargain. Need does not make wrong into good; prosperity does not certify the goodness of the way it was acquired. The poor person's fidelity has a dignity that a comparison of possessions cannot measure.

Thomas explicitly brings Paul's trained contentment into his exposition of the same psalm. He distinguishes spiritual goods, the sufficiency recognized by a contented person, God's provision in necessity and the possibility that providence permits deprivation. His physician analogy concerns God's ordering of the whole person toward salvation.\footnote{Thomas, \work{Super Psalmo} 33, nn. 10--11.} It is not a diagnosis a bystander can make of a particular hungry neighbor. Whether my neighbor should be helped is answered by the command to share, not by speculation about a hidden medical purpose in providence.

\subsection{Paul does not make the gift unnecessary}
Chrysostom's reading of Philippians provides the decisive safeguard. Paul's words about contentment could sound as if the donors have done something superfluous. Paul therefore commends their fellowship in his tribulation. Chrysostom follows the two turns of the argument: the first restrains a benefactor's pride, and the second renews the benefactor's readiness to give. The relationship is preserved because neither donor nor recipient is permitted to dominate it.\footnote{Chrysostom, \work{Homilies on Philippians} XV, ad 4:11--14.}

A recipient can receive help gratefully without treating the donor as the source of salvation. A donor can do real good without treating the gift as a claim to obedience, admiration or control. Paul's security in Christ frees him to recognize the gift's goodness on those terms. His independence is freedom from being owned by circumstances or benefactors, not a refusal of all relationship. The very phrase about sharing his tribulation gives the donor participation in a life that cannot be measured only by what changes hands.

Thomas formulates the objection and answer with unusual clarity. If Paul knows how to endure want, was sending help unnecessary? No: the virtue of enduring poverty does not justify withdrawing assistance.\footnote{Thomas, \work{Super Philippenses} 4, lectio 2.} Contentment is the recipient's virtue, not the neighbor's exemption. This asymmetry matters in Christian speech about suffering. Words a sufferer can truthfully use to describe trust in Christ may become cruel when another person uses them to excuse refusal of help. Paul can say that Christ strengthens him; the Philippians must still decide what to send.

The omitted verses identify the material gift and its bearer, Epaphroditus. They also describe it as an acceptable sacrifice. Chrysostom insists that the money itself cannot buy heavenly things. The donor's merciful purpose and freedom from attachment are what matter; the poor widow's small gift reveals why a monetary scale is inadequate.\footnote{Chrysostom, \work{Philippians} XV, ad 4:15--18.} Thus the value of the gift has two dimensions that must stay together: it supplies a bodily need, and it embodies charity in the giver. To retain only the first makes giving a transaction. To retain only the second risks using the recipient as an instrument of the donor's spiritual advancement.

\subsection{Abundance is also a trial}
Paul has learned both hunger and plenty. Chrysostom refuses the assumption that prosperity requires no discipline. Plenty can make a person idle, boastful or indifferent; want can tempt a person to evil. Paul remains ordered toward his calling in both conditions because his strength comes from Christ.\footnote{Chrysostom, \work{Philippians} XV, ad 4:12--13.} The test is not whether a Christian prefers starvation to food. It is whether the presence or absence of possessions becomes master of the will.

This distinction brings the Gospel's ordinary distractions into sharper focus. Matthew names a farm and business. Neither is condemned simply for existing. The guests make light of the invitation and turn away; ordinary responsibilities become the occasion of refusing the good offered. Gregory describes excessive immersion in earthly labor and longing for gain, while Chrysostom stresses negligent disregard behind the excuses.\footnote{Gregory, \work{Homiliae in Evangelia} 38.5; Chrysostom, \work{Matthew} LXIX.1. Chrysostom also brings in details from Luke which are not in Matthew's appointed text.} A full schedule can conceal spiritual poverty as effectively as a full purse.

The feast does not ask the invited to supply what the king lacks. He has prepared it. Their refusal therefore reveals an inability or unwillingness to value the offered relation. The new guests receive what the first would not receive; their gathering discloses the invitation's freedom from the original guests' status. In the long form, the garment then prevents abundance from being reduced to consumption. The appropriate response includes a changed life. In the short form, the open and mixed gathering stands as the ending, with no claim that the omitted inspection has been heard.

\subsection{A table in the valley and a feast beyond death}
Psalm 23 makes present provision and continuing danger visible at once. The table is prepared in the presence of adversaries. Confidence arises from the host's care there, not from a prior evacuation of every enemy. Augustine's account of the Church's journey gives the valley the breadth of mortal life; his Shepherd is already present so that the faithful can afterward be with him.\footnote{Augustine, \work{Enarrationes} 22.3--6.} Thomas likewise distinguishes the instruction and sacramental food of the journey from the final table of the kingdom.\footnote{Thomas, \work{Super Psalmo} 22, nn. 1--2.}

Isaiah's end of death gives the future a measure that present security cannot satisfy. Even the most splendid earthly banquet leaves its guests mortal. A feast whose fullness removes death cannot be equated with one fortunate interval of prosperity. Thomas's threefold banquet explains how the Church and the soul participate now in a good whose completion belongs to heaven. Christ's Passion gives the present feast its depth: the Church is nourished by the saving gift of one who did not avoid suffering.\footnote{Thomas, \work{Expositio super Isaiam}, cap. 25, on the mountain and the threefold banquet.}

The distinction also gives grief room to tell the truth. Isaiah promises that tears will be wiped away; the promise would lose its force if tears were already treated as a failure of faith. Hope directs mourning toward a good it cannot produce, while the present community shares the mourner's distress. A person need not pretend that loss has ceased in order to receive nourishment along the way. The table can stand amid enemies because God's presence and the wound of present experience are not mutually exclusive.

This is where the second reading adds something that the banquet images alone might leave unstated. God's provision is not only the food placed before the solitary believer. It also takes the form of other people entering his distress. The Philippians' gift is one concrete way a table is made available within the valley. That connection asks the congregation to become generous participants in provision while waiting for a fullness no human generosity can accomplish.

\subsection{Need, hope and the prayer of the Church}
The Entrance names a need deeper than any shortage of possessions: forgiveness. Augustine grounds hope in a redemption the sinner cannot purchase. The Collect then places the work of endurance and generosity within grace's continuing help. The two prayers do not deny material need; they keep its relief within the larger rescue of the person. A well-fed life can still be captive to sin, and a forgiven person can still need bread.

The acclamation asks for an enlightened heart because hope must learn its object. Chrysostom relates the hope of the calling to the power that raised Christ; Thomas distinguishes present hope from the inheritance yet to be possessed.\footnote{Chrysostom, \work{Ephesians} III, ad 1:15--22; Thomas, \work{Super Ephesios} 1, lectio 6.} An inheritance in God differs from an indefinitely extended sequence of purchases. The request for illumination therefore belongs beside the promise of supply: without it, the word abundance can simply magnify existing desires.

Chrysostom interprets Paul's prayer for the donors in relation to necessities. He imagines people with households, children and limited means; it is fitting to desire sufficiency for people who use their goods generously. He expressly distinguishes meeting need from asking that they become rich. Thomas, without cancelling the obligation of present help, places the complete fulfillment of desire in heavenly glory through Christ.\footnote{Chrysostom, \work{Philippians} XV, ad 4:19--20; Thomas, \work{Super Philippenses} 4, lectio 2.} These emphases answer different parts of the promise. The first gives present bodies their due; the second refuses to make present supply the final end.

The Prayer over the Offerings directs finite gifts and prayers toward heavenly glory. The psalm Communion then concentrates on the good of seeking the Lord, a good Augustine identifies through virtues and the living bread. If the Johannine option is used, the language moves directly to future likeness and sight. Augustine explains waiting as an enlargement of desire: the gift is real, yet the recipient is still being made capable of its fullness.\footnote{Augustine, \work{First Epistle of John} IV.5--7.} The two options illuminate different aspects of the same unfinished journey without becoming a single combined chant.

Finally, the Prayer after Communion asks for the fruit of nourishment in divine participation. The Eucharistic gift concerns the person who will go out into bodily needs, obligations and grief, not a spiritual self detached from them. Charity is one present fruit; resurrection and perfected communion are the final horizon. God's abundance sustains action and hope now precisely because it is greater than what can be completed now.

\begin{fourSenses}
\item[Literal.] Paul knows hunger and commends help; the psalm's table stands amid danger; Isaiah promises an end to death still awaited. The Gospel's prepared feast can be refused, and its long form also judges a guest. John's present childhood awaits manifestation.
\item[Allegorical.] Christ nourishes his body through word, sacrament and the members' mutual care. The Church's present banquet participates in the saving gift of his Passion and moves toward the heavenly feast.
\item[Moral.] Receive help without shame or servility; give without patronage. Practice fidelity in want and generosity in plenty. Another person's contentment never makes neglect permissible, and profitable wrongdoing never becomes the measure of the good.
\item[Anagogical.] Necessities are real goods, but the promised fullness is communion with God, vision, resurrection and the end of death. Hope can endure the interval because its final good exceeds every possession lost along the way.
\end{fourSenses}
